\section{A logarithmic sample lower bound at the critical exponent}
\label{sec:critical-logarithmic-sample-lower}

The preceding comparison has a modest actual-posterior consequence.
It is a logarithmic lower bound in the ambient dimension, not an
exponential lower bound in the parity degree.

\begin{corollary}[Full parity and consistent dictators]
\label{cor:critical-logarithmic-sample-lower}
Let $d\ge2$, $L\ge2$, and suppose
\[
 d^{1/[L(L-1)]}>320000L.
\]
Draw $n$ independent uniform inputs from $\{-1,1\}^d$, labeled by
the full parity $f(x)=\prod_{i=1}^d x_i$. At $\gamma=1$, use the
globally margin-conditioned posterior from the model section.
With probability at least
\[
 1-(1-2^{-n})^{d-1}
\]
over the training sample, its Gibbs error converges to $1/2$ as
$N\to\infty$ at this fixed $d,L,n$. Moreover, writing
$p_N(x)=\Pi_{N,1}\{g:g(x)=f(x)\mid S\}$, on the same sample event,
\[
 \liminf_{N\to\infty}
   \Pr_x\{p_N(x)<2/3\}\ge1/4.
\]
Thus for every fixed $L$ and every fixed $\eta>0$, sample sizes
$n\le(1-\eta)\log_2d$ fail this confidence criterion with probability
tending to one as $d\to\infty$, where the width limit is taken first.
\end{corollary}

\begin{proof}
Write the sample as an $n\times d$ matrix of independent signs, and
let $Y\in\{-1,1\}^n$ be the vector of row parities. Conditional on
$Y$, the first $d-1$ columns remain independent uniform sign vectors:
the final column is uniquely determined by them and $Y$. Each of
these $d-1$ coordinates therefore agrees with all labels with
probability $2^{-n}$, independently. This proves the displayed lower
bound for the event that some positive dictator is training consistent.

On that event the minimum training-consistent critical rate is at most
$J_L(\mathbf1)$, by projective symmetry. Every Boolean function depending
on all $d$ original coordinates has at least $d$ relevant augmented
coordinates. Theorem~\ref{thm:critical-fixed-depth-ordering} gives its
rate at least $d^{1/(L-1)}/64$, whereas
$J_L(\mathbf1)\le5000L d^{1/L}$. The assumed strict inequality gives
a positive rate gap. There are finitely many Boolean functions at
fixed $d$, so the exact posterior-rate principle implies that posterior
mass on functions depending on all original coordinates tends to zero.

A function omitting any one original coordinate has uniform error
exactly $1/2$ on full parity, by averaging over that coordinate.
Consequently the posterior-averaged error tends to $1/2$, or equivalently
$\mathbb E_x p_N(x)\to1/2$. Markov's inequality then yields
\[
 \Pr_x\{p_N(x)<2/3\}
 \ge1-\frac32\mathbb E_x p_N(x),
\]
proving the confidence-error assertion. Finally,
$(1-2^{-n})^{d-1}\le\exp[-(d-1)2^{-n}]\to0$ in the stated sample regime.
\end{proof}

The conclusion concerns posterior sampling and the $2/3$ confidence
criterion. It does not imply an error lower bound for posterior-majority
prediction at threshold $1/2$. No claim uniform in simultaneous growing
dimension and finite width is made.
